% Imporditud: tools/import_eksperimendid.py
% Allikas: Eesti füüsikaolümpiaadi eksperimendi ülesannete
%   kogu (Rael Kalda, 2023), ülesanne Ü79, lahendus L79.
\setAuthor{EFO žürii}
\setRound{lõppvoor}
\setYear{1999}
\setNumber{G E1}
\setDifficulty{3}
\setTopic{geomeetriline optika}
\setVahendid{kumerlääts, joonlaud, millimeetripaber}
\setPunktid{12}
\setAllikas{Eksperimendi kogu Ü79, lahendus lk 68}
\setLahendusseis{toores}

\prob{Koondava läätse suurendus}
Teha kindlaks kuidas oleneb koondava läätse suurendus kaugusest läätse ja eseme vahel ning kaugusest läätse ja silma vahel. Tulemus esitada graafiliselt. Suurendus on suurus, mis näitab, mitu korda on kujutise joonmõõtmed suuremad kui eseme joon- mõõtmed.

\hint

\solu
Idee võrrelda mm-paberi jaotusi, mis paistavad läbi luubi ja selle kõrval (3 p), reaalsed mõõtmistulemused (1 p), tulemus, et suurendus ei sõltu läätse ja silma vahekaugusest (1 p), tulemus, et suurendus suureneb kui läätse kaugus objektist suureneb (1 p), sellel suurenemisel on piir (1 p), suurenduse vähemalt kolme väärtuse mõõtmine (1 p), tulemuse graafiline esitamine (õiged teljed, ühikud-tähised 2 p, mittelineaarne tõusev ja sile joon graafikul 1 p), veahinnang 1 p, järeldus 1 p.
\probend
